\documentclass[10pt,a4paper]{amsart}
\usepackage{amsmath,amssymb,amsthm,graphicx}
% Retained minimal preamble for the Stacks Project chapter sites-cohomology.tex (Cohomology on Sites).
% Verbatim excerpt of the chapter preamble: Xy-pic setup (native preamble lines 14-18) and the native
% theorem declarations with the mathematical macros (native preamble lines 154-217). Copied mechanically;
% nothing added, removed or redefined. The wrapper class and the amsmath/amssymb/amsthm/graphicx base
% packages are supplied by the assembler, so the class/externaldocument part of the native preamble is omitted.
% For commutative diagrams we use Xy-pic
\usepackage[all]{xy}

% We use 2cell for 2-commutative diagrams.
\xyoption{2cell}
% Theorem environments.
%
\theoremstyle{plain}
\newtheorem{theorem}[subsection]{Theorem}
\newtheorem{proposition}[subsection]{Proposition}
\newtheorem{lemma}[subsection]{Lemma}

\theoremstyle{definition}
\newtheorem{definition}[subsection]{Definition}
\newtheorem{example}[subsection]{Example}
\newtheorem{exercise}[subsection]{Exercise}
\newtheorem{situation}[subsection]{Situation}

\theoremstyle{remark}
\newtheorem{remark}[subsection]{Remark}
\newtheorem{remarks}[subsection]{Remarks}

\numberwithin{equation}{subsection}

% Macros
%
\def\lim{\mathop{\mathrm{lim}}\nolimits}
\def\colim{\mathop{\mathrm{colim}}\nolimits}
\def\Spec{\mathop{\mathrm{Spec}}}
\def\Hom{\mathop{\mathrm{Hom}}\nolimits}
\def\Ext{\mathop{\mathrm{Ext}}\nolimits}
\def\SheafHom{\mathop{\mathcal{H}\!\mathit{om}}\nolimits}
\def\SheafExt{\mathop{\mathcal{E}\!\mathit{xt}}\nolimits}
\def\Sch{\mathit{Sch}}
\def\Mor{\mathop{\mathrm{Mor}}\nolimits}
\def\Ob{\mathop{\mathrm{Ob}}\nolimits}
\def\Sh{\mathop{\mathit{Sh}}\nolimits}
\def\NL{\mathop{N\!L}\nolimits}
\def\CH{\mathop{\mathrm{CH}}\nolimits}
\def\proetale{{pro\text{-}\acute{e}tale}}
\def\etale{{\acute{e}tale}}
\def\QCoh{\mathit{QCoh}}
\def\Ker{\mathop{\mathrm{Ker}}}
\def\Im{\mathop{\mathrm{Im}}}
\def\Coker{\mathop{\mathrm{Coker}}}
\def\Coim{\mathop{\mathrm{Coim}}}

% Boxtimes
%
\DeclareMathSymbol{\boxtimes}{\mathbin}{AMSa}{"02}

%
% Macros for moduli stacks/spaces
%
\def\QCohstack{\mathcal{QC}\!\mathit{oh}}
\def\Cohstack{\mathcal{C}\!\mathit{oh}}
\def\Spacesstack{\mathcal{S}\!\mathit{paces}}
\def\Quotfunctor{\mathrm{Quot}}
\def\Hilbfunctor{\mathrm{Hilb}}
\def\Curvesstack{\mathcal{C}\!\mathit{urves}}
\def\Polarizedstack{\mathcal{P}\!\mathit{olarized}}
\def\Complexesstack{\mathcal{C}\!\mathit{omplexes}}
% \Pic is the operator that assigns to X its picard group, usage \Pic(X)
% \Picardstack_{X/B} denotes the Picard stack of X over B
% \Picardfunctor_{X/B} denotes the Picard functor of X over B
\def\Pic{\mathop{\mathrm{Pic}}\nolimits}
\def\Picardstack{\mathcal{P}\!\mathit{ic}}
\def\Picardfunctor{\mathrm{Pic}}
\def\Deformationcategory{\mathcal{D}\!\mathit{ef}}

% --- appended 2026-10-12 (Chief): the project own definitions of the reference/slogan/history
% display environments, copied mechanically from the shared Stacks preamble lines 8-12
% (\usepackage{verbatim} plus the three \newenvironment lines, which the chapter preamble keeps
% verbatim from the project preamble). Nothing added, removed or redefined: these environments
% carry no mathematical content and the project itself typesets them as comments. Without these
% three lines a counted statement whose own text uses them fails with Environment slogan undefined.
\usepackage{verbatim}
\newenvironment{reference}{\comment}{\endcomment}
\newenvironment{slogan}{\comment}{\endcomment}
\newenvironment{history}{\comment}{\endcomment}
\usepackage{cleveref}
\crefname{originalreference}{original reference}{original references}
\Crefname{originalreference}{Original reference}{Original references}
\makeatletter
\newlabel{lemma-characterize-limp}{{07A1}{source}{}{sourcepointer-1}{}}
\newlabel{lemma-characterize-limp@cref}{{[originalreference][0][]07A1}{[1][0][]source}{}{}{}}
\makeatother
\crefname{originalreference}{original reference}{original references}
\Crefname{originalreference}{Original reference}{Original references}
\makeatletter
\newlabel{lemma-direct-image-injective-sheaf}{{072Z}{source}{}{sourcepointer-2}{}}
\newlabel{lemma-direct-image-injective-sheaf@cref}{{[originalreference][0][]072Z}{[1][0][]source}{}{}{}}
\makeatother
\crefname{originalreference}{original reference}{original references}
\Crefname{originalreference}{Original reference}{Original references}
\makeatletter
\newlabel{lemma-pullback-injective-pre-limp}{{0D6X}{source}{}{sourcepointer-3}{}}
\newlabel{lemma-pullback-injective-pre-limp@cref}{{[originalreference][0][]0D6X}{[1][0][]source}{}{}{}}
\makeatother
\crefname{originallemma}{original lemma}{original lemmas}
\Crefname{originallemma}{Original lemma}{Original lemmas}
\makeatletter
\newlabel{sites-lemma-preserve-equalizers}{{00XS}{source}{}{sourcepointer-4}{}}
\newlabel{sites-lemma-preserve-equalizers@cref}{{[originallemma][0][]00XS}{[1][0][]source}{}{}{}}
\makeatother
\title{The pullback of an injective module along a continuous and cocontinuous morphism whose lower shriek commutes with fibre products is totally acyclic}
\author{The Stacks Project Authors}
\date{stacks-project:sites-cohomology.tex:local-snapshot-20261011; GNU FDL 1.2 or later}
\begin{document}
\maketitle
\noindent\textit{Editorial note.} All delivered mathematical text is the author's own, copied from the pinned Stacks Project chapter file sites-cohomology.tex; no statement or proof step was rewritten, reordered or removed.\par\medskip
\noindent\textit{Editorial note.} The counted proof is the author's complete argument: condition (1) of the total acyclicity criterion, the identification of the sections over the iterated fibre products via the commutativity of the lower shriek with fibre products, and the closing appeal to the total acyclicity of injective modules together with the converse of the criterion.\par\medskip
\noindent\textit{Editorial note.} The label lemma-pullback-injective-limp names the counted statement, and the proof environment immediately following it in the source is delivered verbatim as its proof.\par\medskip
\noindent\textit{Editorial note.} The counted statement is carried by the author's own lemma environment in the chapter source.\par\medskip
\noindent\textit{Editorial note.} Pre-open gate: python operations/agent-authoring/tools/precard\_gate.py --source sites-cohomology.tex --label lemma-pullback-injective-limp --start 9470 --end 9508 --local-id 3559 returned RESULT=OK (L3559 free, occupancy FREE, public interval clear). Candidate taken from lane-s4/halfB-priority.json.\par\medskip
\noindent\textit{Editorial note.} The statement contains a footnote recording when the commutativity hypothesis holds; the footnote is delivered verbatim with the statement.\par\medskip
\noindent\textit{Editorial note.} The statement and proof are copied from the pinned chapter file.\par\medskip
\noindent\textit{Allowed standard prerequisites (editorial).}\begin{itemize}
\item Criterion for total acyclicity. Same chapter, the lemma characterizing the sheaves that are totally acyclic in terms of the vanishing of the cohomology of their sections over the fibre products of a surjection. Supplies the criterion applied at the start and its converse applied at the end of the proof.
\item Vanishing of the higher cohomology of the pullback of an injective module. Same chapter, the preliminary lemma showing that condition (1) of the total acyclicity criterion holds for the pullback of an injective module. Supplies condition (1) of the criterion.
\item Injectivity implies total acyclicity. Same chapter, the lemma that an injective sheaf of modules is totally acyclic with respect to the terminal morphism. Supplies, applied to the identity morphism, the total acyclicity of the injective module I used in the final step.
\item The lower shriek commutes with fibre products of sheaves of sets. Sites, the lemma that preservation of finite connected limits by the underlying functor makes the lower shriek commute with fibre products. Supplies the hypothesis of the counted lemma and identifies the sections of the pullback over an iterated fibre product with the sections of I over the corresponding iterated fibre product of lower shrieks.
\end{itemize}\medskip
\ifdefined\hypertarget\hypertarget{sourcepointer-1}{}\fi
\noindent\textit{Source reference.} Original reference 07A1: sources/stacks-sites-cohomology/pinned-original-source.tex, source lines 2145--2145. Prerequisite pointer invoked at proof line(s) 9485; it names a standard fact of the same chapter that the counted argument uses, and every step of the counted proof built on it is retained verbatim in the delivered body.\par\medskip
\ifdefined\hypertarget\hypertarget{sourcepointer-2}{}\fi
\noindent\textit{Source reference.} Original reference 072Z: sources/stacks-sites-cohomology/pinned-original-source.tex, source lines 2212--2212. Prerequisite pointer invoked at proof line(s) 9500; it names a standard fact of the same chapter that the counted argument uses, and every step of the counted proof built on it is retained verbatim in the delivered body.\par\medskip
\ifdefined\hypertarget\hypertarget{sourcepointer-3}{}\fi
\noindent\textit{Source reference.} Original reference 0D6X: sources/stacks-sites-cohomology/pinned-original-source.tex, source lines 9278--9278. Prerequisite pointer invoked at proof line(s) 9486; it names a standard fact of the same chapter that the counted argument uses, and every step of the counted proof built on it is retained verbatim in the delivered body.\par\medskip
\ifdefined\hypertarget\hypertarget{sourcepointer-4}{}\fi
\noindent\textit{Source reference.} Original lemma 00XS: sources/stacks-chow/tags.txt, source lines 1230--1230. Prerequisite pointer invoked at proof line(s) 9480; it names a standard fact of a sibling chapter that the counted argument uses, and every step of the counted proof built on it is retained verbatim in the delivered body.\par\medskip
\begin{lemma}
\label{lemma-pullback-injective-limp}
Let $u : \mathcal{C} \to \mathcal{D}$ be a continuous and cocontinuous
functor of sites. Let $g : \Sh(\mathcal{C}) \to \Sh(\mathcal{D})$
be the corresponding morphism of topoi. Let $\mathcal{O}_\mathcal{D}$
be a sheaf of rings and let $\mathcal{I}$ be an injective
$\mathcal{O}_\mathcal{D}$-module. If
$g_!^{Sh} : \Sh(\mathcal{C}) \to \Sh(\mathcal{D})$
commutes with fibre products\footnote{Holds if $\mathcal{C}$
has finite connected limits and $u$ commutes with them, see
Sites, Lemma \ref{sites-lemma-preserve-equalizers}.}, then
$g^{-1}\mathcal{I}$ is totally acyclic.
\end{lemma}

\begin{proof}
We will use the criterion of Lemma \ref{lemma-characterize-limp}.
Condition (1) holds by Lemma \ref{lemma-pullback-injective-pre-limp}.
Let $K' \to K$ be a surjective map of sheaves of sets on $\mathcal{C}$.
Since $g_!^{Sh}$ is a left adjoint,
we see that $g_!^{Sh}K' \to g_!^{Sh}K$ is surjective.
Observe that
\begin{align*}
H^0(K' \times_K \ldots \times_K K', g^{-1}\mathcal{I}) 
& =
H^0(g_!^{Sh}(K' \times_K \ldots \times_K K'), \mathcal{I}) \\
& =
H^0(g_!^{Sh}K' \times_{g_!^{Sh}K} \ldots \times_{g_!^{Sh}K} g_!^{Sh}K',
\mathcal{I})
\end{align*}
by our assumption on $g_!^{Sh}$. Since $\mathcal{I}$ is an injective module
it is totally acyclic by Lemma \ref{lemma-direct-image-injective-sheaf}
(applied to the identity). Hence we can use the converse of
Lemma \ref{lemma-characterize-limp} to see that the complex
$$
0 \to H^0(K, g^{-1}\mathcal{I}) \to H^0(K', g^{-1}\mathcal{I}) \to
H^0(K' \times_K K', g^{-1}\mathcal{I}) \to \ldots
$$
is exact as desired.
\end{proof}

\end{document}
